% Study sheet: Volume III, Chapter 15, Trajectory Curation.
% Numbers: \val{...} from trajectory_curation_studysheet_numbers.tex, written by trajectory_curation_studysheet_numbers.py.
\documentclass{studysheet}
\input{trajectory_curation_studysheet_numbers}
\sheetvolume{3}
\sheetchapter{15}{Trajectory Curation}
\sheetsources{Numbers from the chapter's mlsysim scenarios (hypothetical pass rates, slot times and token counts); assumed inputs are marked.}

\begin{document}
% ===================================================================== FRONT: LEARN
\frontside
\begin{bigidea}
A record says what happened; a training example claims part of it is worth reproducing for a named
learner. Before any gradient, curation decides whether a failure belongs in the weights, which tokens are
\hl{targets} and which \hl{context}, and keeps \hl{evaluation evidence} out of selection. The trained policy keeps the same grants.
\end{bigidea}

\begin{leftcol}
\sheetsec{Key ideas}
\begin{ideas}
\idea{Weights are the last rung.}{Probe context, tool schema, then runtime; a failure vanishing
  under a probe is fixed there. Train only the residual; checks stay.}
\idea{One record, three artifacts.}{Imitation material (selected targets), a preference comparison (shared
  context, justified order), or a rollout fixture (resettable task, grader).}
\idea{A split protects a claim.}{Every attempt, variant and derived record joins one task
  group; whole groups are assigned before collection.}
\idea{Admission is per route.}{Record, use, process, outcome and target checks must all pass for that
  consumer; pending or unknown is never a pass.}
\idea{A transformed record needs a new decision.}{Redaction makes a new version in the same split group,
  pending until revalidated; a valid hash proves unmodified bytes, not truth or permission.}
\idea{Self-selection narrows coverage.}{Kept successes dominate the next round; track action diversity and
  held-out transfer; keep an external anchor.}
\end{ideas}

\vspace{8pt}
\fcard{Cascade cost and order}%
  {\ExpectedCheckingTime = \CheckSlotTime_1 + \CheckPassProbability_1\CheckSlotTime_2
   + \CheckPassProbability_1\CheckPassProbability_2\CheckSlotTime_3 + \cdots}%
  {Check $\AdmissionCheckIndex$ takes $\CheckSlotTime_\AdmissionCheckIndex$ seconds when reached and passes with
   $\CheckPassProbability_\AdmissionCheckIndex$; evaluation stops at the first failure. Yield $\AdmissionRouteYield$ is the
   product of all pass rates; slot time per admitted record is $\ExpectedCheckingTime/\AdmissionRouteYield$.}%
  {Order independent checks by ascending cost per rejection, \mbox{$\CheckSlotTime_\AdmissionCheckIndex/(1-\CheckPassProbability_\AdmissionCheckIndex)$}.
   Adjacent checks cost $\FirstAdjacentCheckSlotTime + \FirstAdjacentCheckPassProbability\SecondAdjacentCheckSlotTime$ in one order,
   $\SecondAdjacentCheckSlotTime + \SecondAdjacentCheckPassProbability\FirstAdjacentCheckSlotTime$ in the other; the admitted set stays.}{}
\end{leftcol}%
\begin{rightcol}
\sheetsec{A recovery record: context and targets}
{\centering
\begin{tikzpicture}[sheet, refused/.style={nbox, dashed}]
  \foreach \i/\e/\t/\sty in {0/e1/{task,\\context}/nbox, 1/e2/{edit\\refused}/refused, 2/e3/{ask for\\schema}/hotbox,
                            3/e4/{fixed\\edit}/hotbox, 4/e5/{check\\summary}/nbox, 5/e6/{report,\\scope}/hotbox}{
    \node[\sty, minimum width=36pt, minimum height=26pt, align=center] (n\i) at ({20+\i*41},30) {\textbf{\e}\\\t};
  }
  \foreach \i in {2,3,5}{\fill[accent] ({20+\i*41-18},6) rectangle ({20+\i*41+18},10);}
  \foreach \i in {0,1,4}{\draw[muted] ({20+\i*41-18},6) rectangle ({20+\i*41+18},10);}
  \node[note, anchor=north] at (122,1) {solid strip: tokens in the loss; the refused edit is context only};
\end{tikzpicture}\par}

\vspace{8pt}
\fcard{Random trace split}%
  {\Pr(\text{straddle}) = 1-\TrainingSplitFraction^{\IndependentAttemptCount}-(1-\TrainingSplitFraction)^{\IndependentAttemptCount}}%
  {$\IndependentAttemptCount$ attempts per task, each sent to training with probability $\TrainingSplitFraction$:
   the share of tasks with attempts on both sides. Their score can reward familiarity, not transfer.}%
  {}%
  {\begin{tikzpicture}[sheet]
     \foreach \yv in {0,50,100}{\draw[gridline] (30,{\yv*0.45}) -- (205,{\yv*0.45});
       \node[small, anchor=east] at (26,{\yv*0.45}) {\yv\%};}
     \foreach \kk in {1,2,4,6,8}{\node[small, anchor=north] at ({30+(\kk-1)*25},-3) {\kk};}
     \draw[accent, line width=1.1pt] plot[domain=1:8, samples=50, smooth]
       ({30+(\x-1)*25},{45*(1-\val{chf}^\x-(1-\val{chf})^\x)});
     \filldraw[fill=accent, draw=accent] (205,{0.45*\val{chstraddle}}) circle (2pt);
     \node[lab, anchor=south east] at (201,{0.45*\val{chstraddle}+2}) {\val{chstraddle}\%};
     \node[note, anchor=north] at (117,-13) {attempts per task $\IndependentAttemptCount$, at $\TrainingSplitFraction = \val{chf}$};
   \end{tikzpicture}}
\fcard{Harvest and contrast}%
  {1-(1-\AttemptPassProbability)^{\IndependentAttemptCount}, \qquad
   \InformativeTaskFraction = 1-(1-\AttemptPassProbability)^{\IndependentAttemptCount}-\AttemptPassProbability^{\IndependentAttemptCount}}%
  {$\AttemptPassProbability$: pass rate per independent attempt. Left: the fixture yields a success to imitate.
   Right: both outcomes, the contrast preference and group-relative training need.}%
  {Mid-difficulty fixtures maximize contrast; easy ones are cheap imitation anchors.}{}
\fcard{Three denominators}%
  {\StoredTokenVolume = \CategoryRecordCount\,\MeanStoredTokensPerRecord, \qquad
   \TargetTokenVolume = \CategoryRecordCount\,\MeanTargetTokensPerRecord}%
  {Per class: $\CategoryRecordCount$ records with $\MeanStoredTokensPerRecord$ stored and $\MeanTargetTokensPerRecord$
   selected target tokens each.}%
  {Specify a mixture in target tokens; budget training compute in admitted stored tokens, storage in all.}{}
\end{rightcol}

% ===================================================================== BACK: PRACTICE
\backside
\begin{answers}
\ans{Finish}{\val{vrejproc} s exceeds \val{rejout} s: outcome first. $\val{reach} \times (\val{vproct} \times \val{outmiss}
  - \val{poutt} \times \val{vprocmiss}) = \val{vgain}$ s saved: \val{veproc} to \val{veout} s (\val{vsavepct}\% less).
  $\val{veout} \div \val{vyield} \approx \val{vadm}$ s per admitted record (cascade card).}
\ans{Random split}{(a) One minus \val{qfkone} minus \val{qmkone}: \val{qsone}\%; one minus \val{qfktwo} (other term
  negligible): \val{qsktwo}\%. (b) Familiarity with siblings, not transfer to unseen tasks. (c) Assign whole task groups before
  collection (split card).}
\columnbreak
\ans{Three runs}{(a) Interface rung: fix the adapter; no training. (b) Residual: a rollout fixture recreating the
  race. (c) Recovery demonstration once process and target review approve: refused edit as context, request and fix as
  targets. Not a pair: no shared starting context. (d) No: a new version, pending until revalidated (key ideas).}
\ans{Harvest}{(a) One minus \val{qmidpow}: \val{qymid}\%; one minus \val{qeasymisspow}: \val{qyeasy}\%. (b) One minus
  \val{qmidpow} twice: \val{qmmid}\%; one minus \val{qeasymisspow} minus \val{qeasypow}: \val{qmeasy}\%. Mid: group-relative
  contrast; easy: cheap imitation anchor (harvest card).}
\ans{Denominators}{Records \val{qnr} of \val{qntot}: \val{qrecshare}\%. Stored \val{qvr} of \val{qvtot} million:
  \val{qstoredshare}\%. Targets \val{qtr} of \val{qttot} million: \val{qtargetshare}\%. Mixture in targets; compute in
  admitted stored tokens (denominators card).}
\end{answers}

\begin{leftcol}
\sheetsec{Worked example}
\begin{practicebox}\boxtitle{Process audit or outcome tests first?}
Only the process audit and outcome tests may swap.\par\vspace{3pt}
\begin{sheettable}{@{}L c c c c c@{}}
\thead{check} & \thead{record} & \thead{use} & \thead{process} & \thead{outcome} & \thead{target}\\ \midrule
pass & \val{precp}\% & \val{pusep}\% & \val{pprocp}\% & \val{poutp}\% & \val{ptargp}\%\\
slot time & \val{prect} s & \val{puset} s & \val{pproct} s & \val{poutt} s & \val{ptargt} s\\
\end{sheettable}
\vspace{4pt}
\begin{steps}
\item Cost per rejection: process \val{pproct} $\div$ \val{procmiss} = \val{rejproc} s; outcome
  \val{poutt} $\div$ \val{outmiss} = \val{rejout} s. \res{Process first.}
\item The pair is entered with \val{precp}\% $\times$ \val{pusep}\% = \val{reach}; only its terms change:
  \val{reach} $\times$ (\val{poutt} $\times$ \val{procmiss} $-$ \val{pproct} $\times$ \val{outmiss}) = \res{\val{gain}~s} saved.
\item Per attempted record \val{eout} s becomes \val{eproc} s (\res{\val{savepct}\% less}). Yield \val{routeyield}
  either way: \val{eproc} $\div$ \val{routeyield} = \res{\val{admproc}~s} per admitted record.
\end{steps}
\takeaway{Order buys speed, not a new corpus.}
\end{practicebox}

\vspace{10pt}
\sheetsec{Your turn: finish it}
\begin{practicebox}\boxtitle{A slower, less selective audit}
Variant: the process audit now takes \val{vproct} s and passes \val{vprocp}\%.\par\vspace{4pt}
\textbf{Given:} its cost per rejection is \val{vproct} $\div$ \val{vprocmiss} = \val{vrejproc} s.\par
\textbf{You:} choose the order; find the saving per attempted record and the time per admitted record.\par\vspace{30pt}
\end{practicebox}
\end{leftcol}%
\begin{rightcol}
\sheetsec{Common mistakes}
\mistake{``The run was accepted, so every token in it is a good target.''}%
  {Acceptance judges the outcome, not each step. The chapter's recovery record stores \val{recstored} tokens, of
   which \val{rectarget} are targets; supervising the refused edit teaches the model to repeat it.}{}
\mistake{``Fine-tuning on every failure costs only \$\val{costall}, so skip the probes.''}%
  {Probes place \val{failsystem} of \val{failall} failures in the tool schema and sandbox pool. Training leaves
   both broken. Train on the \val{failresidual}-failure residual (\$\val{costres}).}{}
\mistake{``Dropping target review is the cheapest speedup.''}%
  {Time per admitted record falls from \val{admproc} to \val{dropadm} s, but yield rises to \val{dropyield}
   because unreviewed targets pass: a different corpus, not a saving.}{}

\vspace{2pt}
\sheetsec{Check yourself}
\begin{qs}
\item \assumed{} each trace goes to training with probability \val{qf}. (a) Share of tasks straddling the split
  at \val{qkone} attempts per task? At \val{qktwo}? (b) What does the test score measure? (c) The fix?
\item Three runs: (a) fails until the search adapter adds pagination markers; (b) survives every probe; no attempt
  solves its race; (c) a refused edit, a schema request, a fix, accepted. Remedy and artifact for
  each? Can (c)'s refused edit and fix form a pair? (d) May (c), with paths redacted, enter training as is?
\item \assumed{} \val{qkthree} attempts per fixture, pass rates \val{qpmid} and \val{qpeasy}. (a) Chance of a success?
  (b) Of a mixed group? Which fixture serves which learner?
\item \assumed{} \val{qnd} direct records (\val{qsd} stored, \val{qwd} target tokens each) and \val{qnr} recovery
  records (\val{qsr} stored, \val{qwr} target). Recovery's share of records, stored tokens, targets? Which sets
  the mixture? The compute?
\end{qs}
\end{rightcol}
\end{document}
